{\footnotesize
\begin{longtable}{>{\raggedright\arraybackslash}p{0.11\textwidth}r>{\raggedright\arraybackslash}p{0.11\textwidth}l>{\raggedright\arraybackslash}p{0.12\textwidth}l>{\raggedright\arraybackslash}p{0.24\textwidth}}
\caption{Uniform-shift scenarios for the western lowlands. Under a spatially uniform change, the smallest cap among sites binds. Cold: smallest downward cap graded \emph{ok}. Warm: smallest upward cap not graded GEO.}\label{tab:scenarios}\\
\toprule
\textbf{Scenario} & \textbf{MTWA shift} & \textbf{Binding site} & \textbf{Env.} & \textbf{Binding taxon} & \textbf{Grade} & \textbf{All western caps (smallest per site)} \\
\midrule
\endfirsthead
\toprule
\textbf{Scenario} & \textbf{MTWA shift} & \textbf{Binding site} & \textbf{Env.} & \textbf{Binding taxon} & \textbf{Grade} & \textbf{All western caps (smallest per site)} \\
\midrule
\endhead
\bottomrule
\endfoot
Coldest permitted & $-$2.50 & Battle Ground Lake & regional & \emph{Tsuga mertensiana} & ok & Battle Ground Lake 2.5, Lake Washington 3.1 \\
Warmest permitted & $+$2.46 & Lake Washington & regional & \emph{Pseudotsuga menziesii} & provisional & Lake Washington 2.5, Hoh valley 4.0, Pender Is.\ (BC) 4.6, Battle Ground Lake 4.9, Killebrew Lake Fen 5.0 \\
\end{longtable}
}
